% Cote 119: offprints, reports, courses, CV and letters (typed and printed part): archive cote 119, site page 201 (PDF page 202).
% typescript. Draft compared: none.
% AI-transcribed by Claude Opus 5.5 (claude-opus-5-5, adaptive thinking, effort high), Message Batches API, batch r119_full (msgbatch_01DwZNm5Y7YsKBWyAN4sFs8w), 2026-11-13; not yet checked by eye.
\scanpage{119/201}
\mspage{- 6 -}

e) \uline{Espaces rigide-analytiques}. M'inspirant de l'exemple de la ``courbe ellip-%
tique de Tate'', et des besoins de la ``géométrie formelle'' sur un anneau de valua-%
tion discrète complet, j'étais parvenu à une formulation partielle de la notion
de variété rigide-analytique sur un corps valué complet, qui a joué son rôle dans
la première étude systématique de cette notion par J. Tate [\unsure{*}]. Par ailleurs,
les ``cristaux'' que j'introduis sur les variétés algébriques sur un corps de
caractéristique $>0$ peuvent s'interpréter parfois en termes de fibrés vecto-%
riels à connexion intégrable sur certains types d'espaces rigide-analytiques
sur des corps de caractéristique nulle; ceci fait pressentir l'existence de
relations profondes entre cohomologie cristalline en car. $>0$ , et cohomologie
de systèmes locaux sur des variétés rigide-analytiques en car. nulle.

\uline{Références} : J.P. Serre, H. Grauert, H. Cartan, P. Deligne, A. Douady, B. Mal-%
grance, K.Knorr, R. Kiehl, J. Tate.

6. \uline{Groupes algébriques} ([16 SGA 3 - en trois volumes] [12 bis])

Ce sujet relève à la fois de la géométrie algébrique et de la théorie
des groupes. Le travail cité SGA 3 se place surtout sur des schémas de base
généraux, et la part de la géométrie algébrique y est certes considérablement
plus large que celle de la théorie des groupes. Néanmoins, grâce à la technique
des schémas, nous y obtenons des résultats nouveaux même dans le cas de groupes
définis sur un \uline{corps} de base, les plus intéressants (relatifs surtout au cas
d'un corps de base imparfait) étant contenus dans SGA 3, Exp XIV. Ma contribu-%
tion principale, continuant dans la voie ouverte par A. Borel et C. Chevalley
dans le contexte de la géométrie algébrique habituelle, a été de montrer le
parti qu'on pouvait tirer d'une application systématique de la théorie des
schémas aux groupes algébriques et aux schémas en groupes.

\uline{Références} : J. Tits, F. Bruhat, M. Demazure, P. Gabriel, A. Borel, D. Mumford.

7. \uline{Groupes discrets} ([18, Exp VIII], [13 bis])

Dans [18, Exp. VIII] je développe une théorie purement algébrique des
classes de Chern desreprésentations d'un groupe discret sur un corps de base
(ou même un anneau de base) quelconque, avec des applications de nature arith-%
métique sur l'ordre des classes de Chern des représentations complexes. Cette
théorie peut être considérée comme cas particulier d'une théorie des classes
de Chern des représentations linéaires de schémas en groupes quelconques,

\margin{201}
